\subsection{强奸创伤综合征与报案决策：受害者会经历什么}

前文已讨论性侵犯的法律定义与心理康复方法（参见《性犯罪与法律责任》《性创伤与心理康复》等章节）。本节补上受害者最需要、但现有内容较少涉及的一块：\textbf{事后会发生什么、要不要报案、以及怎么在这个过程里保住自己的选择权}。

\textbf{一、强奸创伤综合征（Rape Trauma Syndrome）}
1974 年由 Burgess 与 Holmstrom 提出，描述性侵受害者在事后的典型反应模式。它不是一个诊断标签，而是一张"你会经历什么"的地图——知道自己在正常范围内，本身就是一种安慰。
\begin{itemize}
\item \textbf{急性期（事后数小时至数周）}：情绪上表现为恐惧、羞耻、麻木、愤怒交织，也可能出现"异常平静"——\textbf{没有崩溃不代表没受伤}；身体上有失眠、噩梦、食欲改变、肌肉紧张、对特定场景或气味的强烈回避；
\item \textbf{重组期（数周至数月）}：生活逐步恢复，但可能伴随对亲密接触的回避、性欲变化、信任困难、注意力与工作能力下降。这一段最容易出现"我是不是有问题"的自我怀疑；
\item \textbf{长期影响}：多数人能够恢复；但部分人会出现创伤后应激障碍（PTSD）、抑郁或焦虑，需要专业治疗。\textbf{越早获得支持，恢复通常越顺利}——不论你选择是否报案。
\end{itemize}

\textbf{二、报案决策：这不是一道道德题}
"应该报案"和"不报案就是纵容"的说法，把复杂的现实简化成了对受害者的又一次要求。可以做的是把选项和代价摊开：
\begin{itemize}
\item \textbf{选择报案的理由}：证据保全最有效（尤其是现场与物证）、对加害者的法律追责、可能获得法律援助与医疗支持；
\item \textbf{选择不报案的现实考量}：担心隐私与身份暴露、担心被质疑或不被相信、担心家庭压力、举证困难与漫长的程序。这些顾虑都是真实的，不是"软弱"；
\item \textbf{不必二选一}：可以先做\textbf{医学处置与证据保全}（就医取证、保留物证），把"是否报案"的决定往后推——\textbf{取证≠报案}，但不取证往往会永久失去取证的窗口；
\item \textbf{时间窗意识}：72 小时内的医学处置（HIV 阻断 PEP、紧急避孕、性传播感染预防性治疗）与物证采集最为关键；物证通常建议在 72 小时内（越早越好）完成采集。
\end{itemize}

\textbf{三、证据与医疗：怎么做才不留遗憾}
\begin{itemize}
\item \textbf{尽量不洗澡、不刷牙、不更换衣物、不冲洗下身}——这不是"讲究"，是因为很多痕量证据只存在于身体、衣物与洗漱痕迹中；
\item \textbf{保留现场相关物品}：床单、衣物、使用的杯子；威胁或骚扰的聊天记录、通话记录、截图原件（不要只保存转发后的版本）；
\item \textbf{尽快到可做取证的医疗机构}：说明情况，要求进行性侵受害者的医学检查与取证（许多城市设有专门的取证机构或指定医院，可先电话确认）；同时评估 PEP、紧急避孕与性传播感染预防治疗；
\item \textbf{如果已经洗过澡}：仍然可以就医——其他形式的证据（伤痕、病史、对方留下的痕迹）依然有意义，不要因为"证据没了"就放弃求助。
\end{itemize}

\textbf{四、二次创伤：它来自哪里}
\begin{itemize}
\item 常见来源：反复询问细节、被质疑"为什么喝那么多""为什么穿那样""为什么不反抗""为什么现在才说"、亲友的指责或过度渲染，以及把受害者信息在社群中扩散；
\item \textbf{给亲友的规则}：不问细节、不做评判、不代为决定、不传播——先提供安全感和陪伴，把"接下来怎么办"的决定权留给当事人；
\item \textbf{给当事人的提醒}：任何让你感到更糟的交谈或处理方式，都可以拒绝。你有权换医院、换医生、换咨询师、换律师，也有权在任何一个阶段停下。
\end{itemize}

\noindent\textit{【编者按】}本书不鼓励也不责备任何一种选择。"报案或不报案"通常是一次又一次重做的决定，而不是一次性的道德判断；本节想提供的只是——\textbf{在做决定之前，尽可能多的可用信息与可撤回的空间}。
